\section{Conclusions}
\label{sec:conclusions}

Scored as diagnostic tests on 6,663 experimental spectra with known structures, the classical infrared correlation rules divide into two regimes. At or above 1500~\wn{} most rules carry real evidence, with a median positive likelihood ratio of about 3 and values up to 43 for the saturated ester carbonyl; at or below 1350~\wn{} band position alone is close to uninformative, with a median ratio of 1.1 to 1.2, because almost every spectrum has a band in almost every window. Two collections that share no compound, instrument, sampling technique or decade rank the rules the same way (Spearman 0.79; 87\% within a factor of two). Intensity qualifiers raise specificity but cost more sensitivity than the tables acknowledge, and for the nitrile stretch the word ``intense'' turns the most conspicuous band in the spectrum into a rule that recognises one nitrile in ten. Several rules, including the benzene substitution patterns, are of little use for confirming a group but good for excluding it, a direction the tables do not list. The quantified table, the frozen rule set and the scripts are released so that the audit can be extended to other sources and other collections.
